\input{../../stock/en-tete_v5.tex}

\newcounter{chapitre}
\setcounter{chapitre}{2}

\title{\Large Chapitre 2 \\ \Huge Vecteurs aléatoires}

\begin{document}
\input{../../stock/commands.tex}
\maketitle


\newcommand{\tribu}[0]{\mathcal{T}}
\newcommand{\univ}[0]{\Omega}
\newcommand{\proba}[0]{\mathbb{P}}
\newcommand{\probaset}[1]{\mathbb{P}\left(\left\{#1\right\}\right)}



\section{Introduction}
\begin{definition}{}{variable aléatoire}
    Soit $(\univ, \tribu, \proba)$ un espace de probabilité. On appelle \notion{vecteur aléatoire} sur $\univ$ une application de $\univ$ dans $\R^n$.
    Pour $n=1$, on parle de \notion{variable aléatoire}.
\end{definition}

\begin{remarque}{}{notation associée}
    Pour $A \in \R^n$, on notera $\{X \in A\}$ l'évènement~:
    $$\{\omega \in \univ,\, X(\omega) \in A\}$$
\end{remarque}

\begin{definition}{}{vecteur aléatoire discret}
    Soit $(\univ, \tribu, \proba)$ un espace de probabilité. Un vecteur aléatoire $X$ sur $\Omega$ est dit \notion{discret} s'il existe $F \subset \R^n$ au plus dénombrable tel que~:
    $$\proba\Bigl(\{X \in F\}\Bigr) = 1$$
    
\end{definition}


Nous émettrons dans le cadre du cours l'hypothèse suivante.
\begin{remarque}{}{en lien avec la définition}
    Si $X$ est un vecteur aléatoire sur $(\univ, \tribu, \proba)$, pour tout ouvert $A$ de $\R^n$~:
    $$\{\omega \in \univ,\, X(\omega) \in A\} = \{X \in A\} \in \tribu$$
\end{remarque}

\begin{definition}{}{vecteur aléatoire à densité}
    $X$ est un vecteur aléatoire définie sur $(\univ,\, \tribu,\, \proba)$ à valeurs dans $F$ s'il il existe $p : F \to \R_+$ qui vérifie pour tout $A$ ouvert de $F$~:
    $$\proba\Big(\{X \in A\}\Big) = \int_A p(x)\d{x}$$
\end{definition}

\begin{remarque}{}{à ce propos}
    Pour $A = F$,\, la probabilité précédemment évoquée vaut~:
    $$\probaset{X \in F} = 1$$
\end{remarque}
\section{Lois usuelles}

\begin{definition}{}{loi d'une variable aléatoire}
    La loi d'une variable aléatoire discrète $X$ à valeurs dans $F$ (\ie $\proba\Big(\{X \in F\}\Big)=1$) est entièrement définie par la famille de réels positifs~:
    $$\Bigg(\probaset{X = x}\Bigg)_{x \in F}$$
    de somme 1.
\end{definition}

\begin{exemple}{}{loi de Bernoulli}
    Pour une loi de Bernoulli~:
    \begin{itemize}
        \item $\displaystyle F = \{0, 1\}$
        \item $\displaystyle \probaset{X = 1} = p$ et $\probaset{X = 0} = 1-p$
    \end{itemize}
    Une telle variable aléatoire mesure la probabilité de succès d'une épreuve de Bernoulli (succès de probabilité $p$, échec de probabilité $1-p$).
\end{exemple}

\begin{exemple}{}{loi binomiale}
    Pour une loi binomiale~:
    \begin{itemize}
        \item $\displaystyle F = \intint{0}{n}$ où $n \in \N^*$
        \item $\displaystyle \forall k \in \intint{0}{n},\, \probaset{X = k} = \binom{n}{k}p^k(1-p)^{n-k}$
    \end{itemize}
    Une telle variable aléatoire mesure le nombre de succès au terme de la répétition indépendante de $n$ épreuves de Bernoulli du même paramètre $p$.
\end{exemple}

\begin{exemple}{}{loi géométrique}
    Pour une loi géométrique~:
    \begin{itemize}
        \item $\displaystyle F = \N^*$.
        \item $\displaystyle \forall k \in \N^*,\, \probaset{X = k} = p (1-p)^{k-1}$
    \end{itemize}
    En numérotant à partir de 1 une suite d'épreuves de Bernoulli répétées indépendamment et indéfiniment, une telle variable aléatoire mesure l'indice du premier succès. 
\end{exemple}

\begin{exemple}{}{loi de Poisson}
    Pour une loi de Poisson~:
    \begin{itemize}
        \item $\displaystyle F = \N$
        \item $\displaystyle \forall k \in \N,\, \probaset{X = k} = \frac{\lambda^k}{k!}e^{-\lambda}$
        
    \end{itemize}
    Une loi de Poisson a des intérêts multiples, comme celui d'approximer sous certaines conditions une loi binomiale.
\end{exemple}

\begin{remarque}{}{sur ces lois}
    On retrouve par dénombrement l'expression de telles lois. Par exemple dans le cas de la loi géométrique, $\{X = k\}$ correspond à l'évènement "La $k$-ème épreuve est la première à réussir.". Cela revient à avoir vu échouer les $k-1$ épreuves précédentes~:
    $$\probaset{X = k} = p (1-p)^{k-1}$$
\end{remarque}


\section{Lois marginales}

\begin{proposition}{}{formule des lois marginales}
    Soit $X$ et $Y$ deux vecteurs aléatoires discrètes sur $(\univ, \tribu, \proba)$, respectivement à valeurs dans $F$ et $G$. On a~:
    \begin{enumeratebf}
        \item $\displaystyle \forall x \in F,\, \probaset{X=x} = \sum_{g \in G}\probaset{Y = g}$
        \item $\displaystyle \forall g \in G,\, \probaset{Y=g} = \sum_{x \in F}\probaset{X = x}$
    \end{enumeratebf}
\end{proposition}

\begin{remarque}{}{démonstration}
    On doit ceci à la formule des probabilités totales appliquée aux systèmes complets d'évènements $\Big( Y=g \Big)_{g \in G}$ et $\Big( X=x \Big)_{x \in F}$.
\end{remarque}

\section{Variables aléatoires à densité}

Pour une \notion{variable aléatoire à densité}, on dira que la densité caractérise la loi de cette variable aléatoire~:

\begin{definition}{}{variable aléatoire à densité}
    Soit $X$ est une variable aléatoire sur $(\univ, \tribu, \proba)$ à valeurs dans $F$.
    $X$ est qualifiée de \notion{variable aléatoire à densité} s'il existe $p : F \to \seg{0}{1}$, dite \notion{densité de probabilité}, vérifiant pour toute partie $A$ de $F$~:
    $$\probaset{X \in A} = \int_A p(t)\d{t}$$
\end{definition}

\begin{remarque}{}{sur cette définition}
    La donnée de $X$ est alors entièrement caractérisée par sa densité de probabilité $p$.
\end{remarque}

\begin{exemple}{}{loi de densité uniforme}
    Pour une loi uniforme de paramètres $a<b$ dans $\R$~:
    \begin{itemize}
        \item $\displaystyle F = \seg{a}{b}$
        \item $\displaystyle \forall x \in \seg{a}{b},\, \probaset{X = x} = \frac{\mathds{1}_{\seg{a}{b}}(x)}{b-a}$
    \end{itemize}
    La fonction \code{rand} en \code{C} suit une loi uniforme sur $\seg{0}{1}$.
\end{exemple}

\begin{exemple}{}{loi de densité gaussienne}
    Pour une loi gaussienne de paramètres $\mu \in \R$ et $\sigma > 0$~:
    \begin{itemize}
        \item $\displaystyle F = \R$
        \item $\displaystyle \forall x \in \seg{a}{b},\, p(x) = \frac1{\sigma \sqrt{2\pi}}\operatorname \exp\Bigg(-\frac12\left(\frac{x-\mu}{\sigma}\right)^2\Bigg)$
    \end{itemize}
\end{exemple}


\begin{exemple}{}{loi de densité exponentielle}
    Pour une loi exponentielle de paramètre $\lambda > 0$~:
    \begin{itemize}
        \item $\displaystyle F = \R$
        \item $\displaystyle \forall x \in \R,\, p(x) = \lambda e^{-\lambda x} \mathds{1}_{\R_+}(x)$
    \end{itemize}
\end{exemple}

\begin{proposition}{}{formule des lois marginales à densité}
    Soit $(X,Y)$ un couple de variables aléatoires à densité sur $(\univ, \tribu, \proba)$ à valeurs dans $F\times G$. On note $p_{(X,Y)}$ la densité de probablité du couple $(X,Y)$. \\ Alors $p_X$ et $p_Y$ vérifient~:
    \begin{enumeratebf}
        \item $\displaystyle \forall x \in F,\,  p_X(x) = \int_\R p_{(X,Y)}(x,y)\d{y}$
        \item $\displaystyle \forall y \in G,\,  p_Y(y) = \int_\R p_{(X,Y)}(x,y)\d{x}$
    \end{enumeratebf}
\end{proposition}

\begin{demonstration}
    On montre le premier résultat. Montrons que~:
    $$\forall x \in \R,\,  p_X(x) = \int_\R p_{(X,Y)(x,y)\d{y}}$$
    Ce qui revient par définition à montrer que~:
    $$\forall A \subset F,\, \probaset{X \in A} = \int_A \underbrace{\int_\R p_{(X,Y)}(x,y)\d{y}}_{:=p_X(x)} \d{x}$$
    Soit $A \subset F$. On a~:
    \begin{align*}
        \proba(X \in A) &= \proba\Big((X,Y) \in A \times G\Big)\\
        &= \int_{A \times G} p_{(X,Y)}(x,y)\d{(x,y)} \\
        &= \int_A \Bigg(\int_G p_{(X,Y)}(x,y)\d{y}\Bigg)\d{x}
    \end{align*}
    D'où le résultat.
\end{demonstration}


\section{Espérance d'une variable aléatoire}
\subsection{Variable aléatoire discrète}

\begin{definition}{}{famille sommable de réels positifs}
    On dit qu’une \notion{famille $(u_i)_{i \in I}$ de nombres réels positifs est sommable} lorsqu’il existe $M\geq 0$ tel que, pour toute partie finie $J \subset I$ , on ait $\sum_{j \in J}u_j \leq M$.
    On définit alors \notion{la somme de la famille} par~:
    $$\sum_{i \in I}u_i = \sup_{\substack{J \subset I \\ J\text{ finie}}} \sum_{j \in J} u_j$$
\end{definition}

\begin{definition}{}{espérance d'une variable aléatoire discrète}
    Soit $X$ une variable aléatoire discrète à valeurs dans $F \subset \R$.
    $X$ est dite \notion{sommable} si elle admet un moment d'ordre 1 \ie si la famille $\Big(x\probaset{X = x}\Big)_{x \in X}$ est sommable. On notera alors $\mb{E}(X)$ son espérance définie ainsi~:
    $$\mb{E}(X) = \sum_{x \in F} x\probaset{X = x}$$
    
\end{definition}

\begin{exemple}{}{espérance d'une fonction indicatrice}
    Si $A$ est un évènement de $\univ$, alors $\mathds{1}_A : \univ \to \{0, 1\}$ est une variable aléatoire sommable (car $\mathds{1}_A(\univ)$ est fini). Puis~:
    $$\mb{E}(\mathds{1}_A) = \proba(A)$$
\end{exemple}

\begin{theoreme}{}{formule de transfert pour une variable aléatoire discrète}
    Soit $X$ une variable aléatoire discrète sommable à valeurs dans $F \subset \R$ et $g : F \to G$. Alors $g(X)$ est également une variable aléatoire sommable et~:
    $$\mb{E}\Big(g(X)\Big) = \sum_{x \in F}g(x)\probaset{X = x}$$
\end{theoreme}

\subsection{variable aléatoire à densité}

\begin{definition}{}{intégrabilité d'une variable aléatoire à densité}
    Soit $X$ une variable aléatoire à densité à valeurs dans $F$ un espace vectoriel de dimension $n$. On dit de $X$ qu'elle est \notion{intégrable} si l'intégrale~:
    $$\int_F \norme{x} p(x) \d{x}$$
    converge. Le cas échéant on dit que $X$ admet une espérance $\mb{E}(X)$ vérifiant~:
    $$\mb{E}(X) = \underbrace{\int_F \underbrace{x}_{\in F} \underbrace{p(x)}_{\in \R_+}\d{x}}_{\in F}$$
\end{definition}

\begin{remarque}{}{sur cette définition}
    Il est commun d'avoir $F = \R$ muni de la valeur absolue classique ou $F = \R^n$ peu importe la norme (\textit{cf.} dimension finie).
\end{remarque}

\subsection{exemples}


\begin{exemple}{}{espérance d'une loi de Bernoulli}
    Si $X$ suit une loi de Bernoulli de paramètre $p \in \seg{0}{1}$, alors $X(\Omega) = \{0, 1\}$ est fini, $X$ est donc sommable et~:
    \begin{align*}
        \mb{E}(X) &= 0 \times (1-p) + 1 \times p \\
        &= p
    \end{align*}
\end{exemple}

\begin{exemple}{}{espérance d'une loi binomiale}
    Si $X$ suit une loi binomiale de paramètres $n \in \N^*$ et $p \in \seg{0}{1}$, alors $X(\Omega) = \intint{0}{n}$ est fini, $X$ est donc sommable et~:
    \begin{align*}
        \mb{E}(X) &= \sum_{k=1}^n k \binom{n}{k} p^k (1-p)^{n-k} \qquad &\text{premier terme nul}\\
        &= \sum_{k=1}^n \frac{n!}{(k-1)!(n-k)!} p^k (1-p)^{n-k} \\
        &= np \sum_{k=1}^n \frac{(n-1)!}{(k-1)!(n-k)!} p^{k-1} (1-p)^{n-k} \\
        &= np \sum_{j=0}^{n-1} \frac{(n-1)!}{j!(n-j-1)!} p^j (1-p)^{n-1-j}\qquad &\text{changement d'indice}\\
        &= np \Big(p + (1-p)\Big)^{n-1} \qquad &\text{formule du binôme de Newton}\\
        &= np
    \end{align*}
    d'où le résultat.
\end{exemple}

\begin{exemple}{}{espérance d'une loi de Poisson}
    Si $X$ suit une loi de Poisson de paramètres $n \in \N^*$ et $p \in \seg{0}{1}$, alors $X(\Omega) = \N$ est dénombrable, $X$ est donc sommable seulement si la quantité suivante admet une limite finie quand $N$ tend vers $+\infty$~:

    \begin{align*}
        \sum_{k=1}^N k \frac{\lambda^k}{k!}e^{-\lambda}
        &= \sum_{k=1}^N \frac{\lambda^k}{(k-1)!}e^{-\lambda}\\
        &= \lambda e^{-\lambda} \sum_{k=1}^N \frac{\lambda^{k-1}}{(k-1)!}\\
        &= \lambda e^{-\lambda} \sum_{j=0}^{N-1} \frac{\lambda^j}{j!}\\
        &\xrightarrow[N \to +\infty]{} \lambda e^{-\lambda} e^{\lambda} = \lambda
    \end{align*}
    $X$ est donc effectivement sommable, et son espérance vaut~:
    $$\mb{E}(X) = \lambda$$
\end{exemple}

\begin{exemple}{}{espérance d'une loi géométrique}
    Si $X$ suit une loi de géométrique de paramètre $p \in ]0; 1[$ (on évite les cas pathologiques), alors $X(\Omega) = \N$ est dénombrable, $X$ est donc sommable seulement si la quantité suivante admet une limite finie quand $N$ tend vers $+\infty$~:

    \begin{align*}
        \sum_{k=1}^N k p (1-p)^{k-1} = p \sum_{k=1}^N k (1-p)^{k-1}\\
    \end{align*}
    Ceci implique la dérivée d'une somme géométrique de raison $(1-p)$, valant~:
    $$\sum_{k=0}^{+\infty} (1-p)^k = \frac{1}{1-(1-p)} = \frac{1}{p}$$
    Par dérivation terme à terme on a bien sommabilité de $X$, et on trouve~:
    $$\mb{E}(X) = p\sum_{k=1}^{+\infty} k (1-p)^{k-1} = p\frac{1}{p^2} = \frac{1}{p}$$
\end{exemple}

\begin{exemple}{}{espérance d'une loi à densité exponentielle}
    Si $X$ suit une loi exponentielle de paramètre $\lambda > 0$, alors $X(\Omega) = \R_+$ puis~:
    \fonction{f}{\R_+}{\R}{x}{x\lambda e^{-\lambda x}}
    est continue par morceaux, avec par croissances comparées~:
    $$\lim_{x \to \infty} x^2 f(x) = 0$$
    donc~:
    $$f(x) \underset{x \to \infty}{=} o\left(\frac{1}{x^2}\right)$$
    et $X$ est intégrable par comparaison avec une fonction de référence. Son espérance vaut alors~:
    \begin{align*}
        \mb{E}(X) &= \int_0^{+\infty} x \lambda e^{-\lambda x}\d{x}\\
        &= \left[-xe^{-\lambda x}\right]_0^{+\infty} + \int_0^{+\infty} e^{-\lambda x}\d{x} \qquad &\text{intégration par parties}\\
        &= 0 + \left[-\frac{1}{\lambda}e^{-\lambda x}\right]_0^{+\infty}\\
        &= \frac{1}{\lambda}
    \end{align*}
    d'où le résultat.
\end{exemple}

\begin{exemple}{}{espérance d'une loi à densité gaussienne}
    Si $X$ suit une loi gaussienne de paramètres $\mu \in \R$ et $\sigma > 0$, alors $X(\Omega) = \R$ puis~:
    \fonction{f}{\R}{\R}{x}{x \frac1{\sigma \sqrt{2\pi}}\operatorname \exp\Bigg(-\frac12\left(\frac{x-\mu}{\sigma}\right)^2\Bigg)}
    est continue par morceaux, avec par croissances comparées~:
    $$\lim_{x \to \infty} x^2 f(x) = 0$$
    donc~:
    $$f(x) \underset{x \to \infty}{=} o\left(\frac{1}{x^2}\right)$$
    et $X$ est intégrable par comparaison avec une fonction de référence. Son espérance vaut alors~:
    \begin{align*}
        \mb{E}(X) &= \int_{-\infty}^{+\infty} x \frac1{\sigma \sqrt{2\pi}}\operatorname \exp\Bigg(-\frac12\left(\frac{x-\mu}{\sigma}\right)^2\Bigg)\d{x}\\
        &= \int_{-\infty}^{+\infty} (\sigma t + \mu) \frac1{\sigma \sqrt{2\pi}}\operatorname \exp\Bigg(-\frac12 t^2\Bigg) \sigma\d{t} \qquad\text{changement de variable } t = \frac{x-\mu}{\sigma}\\
        &= \int_{-\infty}^{+\infty} t \frac1{\sqrt{2\pi}}\operatorname \exp\Bigg(-\frac12 t^2\Bigg)\d{t} + \mu \int_{-\infty}^{+\infty}  \frac1{\sqrt{2\pi}}\operatorname \exp\Bigg(-\frac12 t^2\Bigg)\d{t}\\
        &= 0 + \mu\int_{-\infty}^{+\infty}  \frac1{\sqrt{2\pi}}\operatorname \exp\Bigg(-\frac12 t^2\Bigg)\d{t} \qquad\text{la première intégrale est nulle par imparité}\\
        &= \mu \qquad \text{la seconde intégrale vaut 1 (on le montre tout de suite après)}
    \end{align*}
    Enfin montrons que $\displaystyle I = \int_{-\infty}^{+\infty}  \frac1{\sqrt{2\pi}}\operatorname \exp\Bigg(-\frac12 t^2\Bigg)\d{t} = 1$. On a~:
    \begin{align*}
        I^2 &= \frac{1}{2\pi} \int_{-\infty}^{+\infty} \int_{-\infty}^{+\infty} \operatorname \exp\Bigg(-\frac12 (x^2 + y^2)\Bigg)\d{x}\d{y} \qquad\text{Fubini pour intégrales}\\
        &= \frac{1}{2\pi} \int_0^{2\pi} \int_0^{+\infty} re^{-\frac12 r^2}\d{r}\d{\theta} \qquad\text{changement de variable en coordonnées polaires}\\
        &= \frac{1}{2\pi} \times 2\pi \times \left[-e^{-\frac12 r^2}\right]_0^{+\infty}\\
        &= 1
    \end{align*}
    d'où $I = 1$ puis $\mb{E}(X) = \mu$.
\end{exemple}

\begin{remarque}{}{sur la démonstration précédente}
    En passant en coordonnées polaires, on a utilisé le fait que~:
    $$\d{(x,y)} = r\d{r}\d{\theta}$$
    Ceci est une application du théorème de changement de variable en dimension quelconque~:
    $$\int_{\varphi(U)} f(x)\d{x} = \int_{U} f\Big(\varphi(x)\Big)\Bigg\vert\det\Big(J_\varphi(x)\Big)\Bigg\vert\d{x}$$
    où ici~: 
    \fonction{\varphi}{U = \R_+ \times [0; 2\pi[}{\R^2}{(r,\theta)}{\Big(r \cos(\theta), r \sin(\theta)\Big)}
    est un $\mc{C}^1$-difféomorphisme, de jacobien $\det\Big(J_\varphi(r,\theta)\Big)$ en $(r, \theta)$ valant $r$, d'où le résultat.
    
\end{remarque}




\section{Indépendance de variables aléatoires}

\begin{definition}{}{indépendance de variables aléatoires}
    Soit $X$ et $Y$ une famille de variables aléatoires sur $(\univ, \tribu, \proba)$ à valeurs dans $F$ et $G$ respectivement. On dit que $X$ et $Y$ sont \notion{indépendantes} si pour tous $x \in F$ et $y \in G$~:
    $$\probaset{(X,Y) = (x,y)} = \probaset{X=x}\probaset{Y = y}$$
    
\end{definition}

\begin{theoreme}{}{caractérisation de l'indépendance de variables aléatoires}
    Soit $X$ et $Y$ une famille de variables aléatoires sur $(\univ, \tribu, \proba)$ à valeurs dans $F$ et $G$ respectivement. Les assertions suivantes sont équivalentes~:
    \begin{enumeratebf}
        \item $X$ et $Y$ sont indépendantes.
        \item Pour toutes fonctions $f$ et $g$ bornées de $F$ et $G$ respectivement dans $\R$~:
        $$\mb{E}\Big(f(X)g(Y)\Big) = \mb{E}\Big(f(X)\Big)\mb{E}\Big(g(Y)\Big)$$
    \end{enumeratebf}
    
\end{theoreme}

\begin{remarque}{}{sur ce théorème}
    Ce théorème sert en particulier dans son sens direct.
\end{remarque}

% 2 : lois à densité

\begin{proposition}{}{caracatérisation de l'indépendance de variables aléatoires à densité}
    Soit $X_1,\, \dots,\, X_n$ des variables aléatoires de densités respectives $p_1,\, \dots,\,  p_n$. Les $X_i$ sont mutuellement indépendantes si et seulement si $(X_1,\, \dots,\, X_n)$ est une variable aléatoire à densité de densité de probabilité~:
    $$p : (x_1,\, \dots,\, x_n) \mapsto \prod_{i=1}^np_i(x_i)$$
\end{proposition}

\section{Propriétés de l'espérance}

\begin{theoreme}{}{linéarité de l'espérance}
    Soit $X$ et $Y$ deux variables aléatoires discrètes sommables (resp. intégrables) sur $(\univ, \tribu, \proba)$ à valeurs dans $F$. Alors pour tous $(\alpha, \beta) \in \R^2$, $\alpha X + \beta Y$ est une variable aléatoire discrète sommable (resp. intégrable) et~:
    $$\mb{E}(\alpha X + \beta Y) = \alpha \mb{E}(X) + \beta \mb{E}(Y)$$
\end{theoreme}

\begin{theoreme}{}{croissance de l'espérance}
    Soit $X$ et $Y$ deux variables aléatoires réelles discrètes sommables (resp. intégrables) sur $(\univ, \tribu, \proba)$. Si $X \leq Y$ presque sûrement, alors~:
    $$\mb{E}(X) \leq \mb{E}(Y)$$
\end{theoreme}

\begin{theoreme}{}{de comparaison d'espérances}
    Soit $X$ et $Y$ deux variables aléatoires sur $(\univ, \tribu, \proba)$ à valeurs dans $F$ telles que~:
    \begin{itemize}
        \item $Y$ est sommable (resp. intégrable)
        \item $\norme{X}_F \leq \norme{Y}_F$ presque sûrement.
    \end{itemize}
    Alors $X$ est sommable.
\end{theoreme}

\subsection{Des inégalités}

\begin{proposition}{}{inégalité de Markov}
    Soit $X$ une variable aléatoire réelle positive sommable (resp. intégrable) sur $(\univ, \tribu, \proba)$. Alors~:
    $$\forall a > 0,\, \probaset{X \geq a} \leq \frac{\mb{E}(X)}{a}$$
\end{proposition}

\begin{demonstration}
    On a l'inégalité~:
    $$a \mathds{1}_{Z \geq a} \leq Z\mathds{1}_{Z \geq a} \leq Z$$
    Par croissance et linéarité de l'espérance~:
    $$a \probaset{\mathds{1}_Z} \leq \mb{E}(Z)$$
    Ainsi on a le résultat~:
    $$ \proba(Z \geq a) \leq \frac{\mb{E}(Z)}{a}$$
\end{demonstration}

\begin{remarque}{}{importante ici}
    Si $F \neq \R$ il n'y a qu'à considérer la variable aléatoire $\norme{X}_F$, bien à valeurs dans $\R_+$ et lui appliquer le résultat. Cette astuce est valable dans tous les énoncés de cette section.
\end{remarque}

\begin{proposition}{}{inégalité de Bienaymé-Tchebychev}
    Soit $X$ une variable aléatoire réelle sommable sur $(\univ, \tribu, \proba)$ admettant une variance finie (ou si $X^2$ est sommable). Alors~:
    $$\forall \varepsilon > 0,\, \probaset{\vert X - \mb{E}(X) \vert \geq \varepsilon} \leq \frac{\mb{V}(X)}{\varepsilon^2}$$
\end{proposition}

\begin{demonstration}
    On applique l'inégalité de Markov à la variable aléatoire $(X - \mb{E}(X))^2$ qui est positive et sommable. On a~:
    $$\probaset{\vert X - \mb{E}(X) \vert \geq \varepsilon} = \probaset{(X - \mb{E}(X))^2 \geq \varepsilon^2} \leq \frac{\mb{E}\Big((X - \mb{E}(X))^2\Big)}{\varepsilon^2} = \frac{\mb{V}(X)}{\varepsilon^2}$$
\end{demonstration}


\begin{theoreme}{}{inégalité de Cauchy-Schwarz}
    Soit $X$ et $Y$ deux variables aléatoires réelles sommables sur $(\univ, \tribu, \proba)$ telles que $X^2$ et $Y^2$ soient sommables. Alors~:
    $$\vert \mb{E}(XY) \vert \leq \sqrt{\esperance{X^2} \esperance{Y^2}}$$
\end{theoreme}

\begin{remarque}{}{cas particulier}
    En appliquant le théorème pour $X$ réelle positive quelconque et $Y = \overline{1}$~:
    $$\mb{E}(X) \leq \sqrt{\mb{E}(X^2)}$$
\end{remarque}

\begin{demonstration}
    On considère pour tout $t \in \R$ la variable aléatoire~:
    $$Z_t = (X + tY)^2 \geq \overline{0}$$
    qui est sommable par linéarité. Par croissance de l'espérance~:
    $$0 \leq \mb{E}(Z_t) = \mb{E}(X^2) + 2t\mb{E}(XY) + t^2\mb{E}(Y^2)$$
    Le polynôme en $t$ de second degré $\mb{E}(Y^2)t^2 + 2\mb{E}(XY)t + \mb{E}(X^2)$ est donc à discriminant négatif ou nul~:
    $$4\mb{E}(XY)^2 - 4\mb{E}(Y^2)\mb{E}(X^2) \leq 0$$
    En effet, si le discriminant était strictement positif, le polynôme s'annulerait en deux points distincts, et serait négatif entre ces deux points. Or on a positivité du polynôme pour tout $t$. \\\\Le résultat découle de cette dernière égalité~:
    \begin{align*}
        &\esperance{XY}^2 \leq \esperance{X^2} \esperance{Y^2}\\
        \text{donc}\quad&\Big\vert\esperance{XY}\Big\vert \leq \sqrt{\esperance{X^2} \esperance{Y^2}}
    \end{align*}
\end{demonstration}

\begin{remarque}{}{condition suffisante pour que $XY$ soit sommable}
    L'inégalité de Cauchy-Schwarz indique clairement que si $X^2$ et $Y^2$ sont sommables, alors il en est de même pour $XY$. Une autre façon de le démontrer est l'inégalité de Young, qui porte directement sur les variables, et non plus sur les espérances~:
    $$\vert XY \vert \leq \frac{X^2}{2} + \frac{Y^2}{2}$$
    Cette dernière inégalité se démontre en sachant que $(X-Y)^2 \geq 0$.
\end{remarque}



\section{Variance et covariance}

\begin{definition}{}{variance d'une variable aléatoire}
    Soit $X$ une variable aléatoire réelle sommable sur $(\univ, \tribu, \proba)$. On dit que $X$ admet \notion{une \notion{variance}} si $X^2$ est sommable. Dans ce cas on note $\mb{V}(X)$ la variance de $X$ définie par~:
    $$\mb{V}(X) = \mb{E}\Big((X - \mb{E}(X))^2\Big)$$
\end{definition}

\begin{proposition}{}{formule de König-Huygens}
    Soit $X$ une variable aléatoire réelle sommable sur $(\univ, \tribu, \proba)$. Si $X$ admet une variance, alors~:
    $$\mb{V}(X) = \mb{E}(X^2) - \mb{E}(X)^2$$
\end{proposition}

\begin{demonstration}
    On développe~:
    \begin{align*}
        \mb{V}(X) &= \mb{E}\Big((X - \mb{E}(X))^2\Big)\\
        &= \mb{E}(X^2) - 2\mb{E}(X)\mb{E}(X) + \mb{E}(X)^2\\
        &= \mb{E}(X^2) - \mb{E}(X)^2
    \end{align*}
\end{demonstration}

\begin{definition}{}{covariance de deux variables aléatoires}
    Soit $X$ et $Y$ deux variables aléatoires réelles sommables sur $(\univ, \tribu, \proba)$ telles que $XY$ est sommable. On appelle \notion{covariance de $X$ et $Y$} la quantité~:
    $$\mr{Cov}(X,Y) = \mb{E}\Big(\big(X - \mb{E}(X)\big)\big(Y - \mb{E}(Y)\big)\Big)$$
    En développant, on trouve~:
    $$\mr{Cov}(X,Y) = \mb{E}(XY) - \mb{E}(X)\mb{E}(Y)$$
    Le résultat est analogue à la formule de König-Huygens.
\end{definition}

\begin{proposition}{}{variance de la somme de deux variables aléatoires}
    Soit $X$ et $Y$ deux variables aléatoires réelles sommables sur $(\univ, \tribu, \proba)$ telles que $X^2$ et $Y^2$ sont sommables. On a~:
    $$\mb{V}(X+Y) = \mb{V}(X) + \mb{V}(Y) + \mr{Cov}(X,Y)$$
    Si $X$ et $Y$ sont indépendantes, alors~:
    $$\mb{V}(X+Y) = \mb{V}(X) + \mb{V}(Y)$$
\end{proposition}

\begin{demonstration}
    On a~:
    \begin{align*}
        \mb{V}(X+Y) &= \mb{E}\Big((X+Y)^2\Big) - \mb{E}(X+Y)^2\\
        &= \mb{E}(X^2) + 2\mb{E}(XY) + \mb{E}(Y^2) - \Big(\mb{E}(X) + \mb{E}(Y)\Big)^2\\
        &= \mb{E}(X^2) + 2\mb{E}(XY) + \mb{E}(Y^2) - \Big(\mb{E}(X)^2 + 2\mb{E}(X)\mb{E}(Y) + \mb{E}(Y)^2\Big)\\
        &= \underbrace{\mb{E}(X^2) - \mb{E}(X)^2}_{=\mb{V}(X)} + \underbrace{\mb{E}(Y^2) - \mb{E}(Y)^2}_{=\mb{V}(Y)} + 2\underbrace{\Big(\mb{E}(XY) - \mb{E}(X)\mb{E}(Y)\Big)}_{=\mr{Cov}(X,Y)}\\
    \end{align*}
    Puis si $X$ et $Y$ sont indépendantes, $\mb{E}(XY) = \mb{E}(X)\mb{E}(Y)$ donc $\mr{Cov}(X,Y) = 0$. D'où les résultats.
\end{demonstration}

\section{Outils pour le calcul de lois de variables aléatoires}

\subsection{fonction génératrice des moments}


Inutile mais je pose ça là.
\begin{definition}{}{moment d'une variable aléatoire}
    Soit $X$ une variable aléatoire à valeurs dans $F$ sur $(\univ, \tribu, \proba)$. On dit que $X$ admet un \notion{moment d'ordre $n$} si $\norme{X^n}_F$ est sommable. Dans ce cas on note $\mb{E}(X^n)$ le moment d'ordre $n$ de $X$. En particulier~:
    \begin{itemize}
        \item être sommable revient à admettre un moment d'ordre 1
        \item admettre une variance revient à admettre un moment d'ordre 2.
    \end{itemize}    
\end{definition}

\begin{definition}{}{fonction génératrice pour une VARD}
    Soit $X$ une variable aléatoire réelle discrète admettant des moments d'ordre $n$ pour tout $n \in \N$. On appelle \notion{fonction génératrice de $X$} la fonction~:
    \fonction{G_X}{\R}{\R}{t}{\mb{E}\Big(t^X\Big)}
    Ainsi car $X$ est discrète~:
    $$G_X(t) = \sum_{x \in X(\omega)} \probaset{X=x} t^x$$
\end{definition}

\begin{exemple}{}{fonction génératrice d'une loi de Bernoulli}
    Si $X$ suit une loi de Bernoulli de paramètre $p \in \seg{0}{1}$, alors~:
    $$G_X(t) = (1-p) + pt$$
\end{exemple}

\begin{exemple}{}{fonction génératrice d'une loi binomiale}
    Si $X$ suit une loi binomiale de paramètres $n \in \N^*$ et $p \in \seg{0}{1}$, alors~:
    \begin{align*}
        G_X(t) &= \sum_{k=0}^n \binom{n}{k} p^k (1-p)^{n-k} t^k \\
        &= \sum_{k=0}^n \binom{n}{k} (pt)^k (1-p)^{n-k} \\
        &= (1-p + pt)^n \qquad &\text{formule du binôme de Newton}
    \end{align*}
\end{exemple}

\begin{proposition}{}{fonction génératrice d'une somme de VARD indépendantes}
    Soit $X$ et $Y$ des variables aléatoires réelles discrètes indépendantes de fonctions génératrices $G_X$ et $G_Y$ respectivement. Alors la fonction génératrice de $X + Y$ est ~:
    $$G_{X+Y} = G_X \times G_Y$$
    
\end{proposition}

\begin{exemple}{}{loi binomiale}
    Soit $X_1,\, \dots,\, X_n$ des variables aléatoires indépendantes et identiquement distribuées (IID) suivant une loi de Bernoulli de paramètre $p \in \seg{0}{1}$. Alors la fonction génératrice de $Y = \sum_{i=1}^n X_i$ est~:
    \begin{align*}
        G_Y(t) &= \prod_{i=1}^n G_{X_i}(t) \qquad &\text{par indépendance}\\
        &= \Big((1-p) + pt\Big)^n \qquad &\text{car les } X_i \text{ suivent une loi de Bernoulli de même paramètre}\\
    \end{align*}
    On retrouve bien le résultat.
\end{exemple}

\begin{exemple}{}{fonction génératrice d'une loi de Poisson}
    Si $X$ suit une loi de Poisson de paramètre $\lambda > 0$, alors~:
    \begin{align*}
        G_X(t) &= \sum_{k=0}^{+\infty} \frac{\lambda^k}{k!}e^{-\lambda} t^k \\
        &= e^{-\lambda} \sum_{k=0}^{+\infty} \frac{(\lambda t)^k}{k!} \\
        &= e^{-\lambda} e^{\lambda t} \qquad &\text{série exponentielle}\\
        &= e^{\lambda(t-1)}
    \end{align*}
\end{exemple}

\begin{exemple}{}{fonction génératrice d'une loi géométrique}
    Si $X$ suit une loi géométrique de paramètre $p \in ]0; 1[$, alors~:
    \begin{align*}
        G_X(t) &= \sum_{k=1}^{+\infty} p(1-p)^{k-1} t^k \\
        &= pt \sum_{k=0}^{+\infty} \Big((1-p)t\Big)^k \\
        &= \frac{pt}{1-(1-p)t} \qquad &\text{série géométrique}\\
    \end{align*}
    On pourra retenir, en notant $q = 1-p$~:
    $$G_X(t) = \frac{pt}{1-qt}$$
    
\end{exemple}

\begin{exemple}{}{}
    On considère $(X_i)_{i \in \N^*}$ une suite de variables aléatoires IID suivant une loi de géométrique de paramètre $p \in \seg{0}{1}$. On note $N$ une variable aléatoire indépendante des $X_i$, à valeurs dans $\N$. On pose alors~:
    $$S : \omega \mapsto \begin{cases*}
    0 & si $N(\omega) = 0$\\
    \displaystyle \sum_{i=1}^{N(\omega)} X_i(\omega) & sinon
    \end{cases*}$$
    Déterminons la fonction génératrice de $S$ en fonction de celle de $N$. Deux démonstrations sont proposées ici.
\end{exemple}

\begin{demonstration}
    Première démonstration proposée par M. DUFOUR, en manipulant les espérances. Par définition~:
    \begin{align*}
        G_S(t) &= \mb{E}(t^S)\\
        &=\mb{E}(\indicatrice_\univ t^S)\\
        &= \esperance{\indicatrice_{\bigsqcup_{n=0}^{+\infty} \{N=n\}} t^S} \qquad &\text{utilisation implicite de la FPT}\\
        &= \esperance{\sum_{n=0}^{+\infty} \indicatrice_{\{N=n\}} t^S}\\
        &= \sum_{n=0}^{+\infty} \esperance{\indicatrice_{\{N=n\}} t^S} \qquad &\text{linéarité de l'espérance}\\
        &= \sum_{n=1}^{+\infty} \esperance{\indicatrice_{\{N=n\}} t^{\sum_{i=1}^N X_i}} + \esperance{\indicatrice_{\{N = 0\}}t^0} \qquad &\text{définition de } S\\
        &= \sum_{n=1}^{+\infty} \esperance{\indicatrice_{\{N=n\}} t^{\sum_{i=1}^n X_i}} + \esperance{\indicatrice_{\{N = 0\}}t^0} \qquad &\text{dans le contexte de la somme, } N=n\\
        &= \sum_{n=1}^{+\infty} \esperance{\indicatrice_{\{N=n\}}} \esperance{  t^{\sum_{i=1}^n X_i}} + \esperance{\indicatrice_{\{N = 0\}}} \qquad &\text{indépendance de } N \text{ avec les } X_i\\
        &= \sum_{n=1}^{+\infty} \probaset{N=n} G_{X_1 + \dots + X_n}(t) + \probaset{N=0} \qquad &\text{esp. de } \indicatrice \text{ et déf. fonction génératrice}\\
        &= \sum_{n=1}^{+\infty} \probaset{N=n} G_{X_1}(t)^n + \probaset{N=0} \qquad &\text{car les } X_i \text{ sont IID}\\
        &= \sum_{n=0}^{+\infty} \probaset{N=n} G_{X_1}(t)^n \qquad &\text{en rajoutant le terme } n=0\\
        &= G_N\Big(G_{X_1}(t)\Big) \qquad &\text{définition de la fonction génératrice}
    \end{align*}
\end{demonstration}

\begin{demonstration}
    Proposition alternative, en manipulant les probabilités. On a $S(\univ) = \N$. Puis par définition~:
    $$G_S(t) = \sum_{n=0}^{+\infty} \proba(S=n)t^n$$
    Pour $n \in \N$, calculons donc $\probaset{S=n}$. On a par formule des probabilités totales appliquée au système complet d'évènements $\Big( \{N = k\}\Big)_{k \in \N}$~:
    \begin{align*}
        \probaset{S=n} &= \sum_{k=0}^{+ \infty}\proba\Big(\{S = n\} \cap \{N = k\}\Big) \qquad\\
        &= \sum_{k=0}^{+ \infty}\proba\left(\left\{\sum_{i=1}^N X_i = n\right\} \cap \{N = k\}\right)\\
        &= \sum_{k=0}^{+ \infty}\proba\left(\left\{\sum_{i=1}^k X_i = n\right\} \cap \{N = k\}\right)\qquad &\text{car dans la probabilité, } N = n\\
        &= \sum_{k=0}^{+ \infty}\proba\left(\left\{\sum_{i=1}^k X_i = n\right\}\right) \proba\left(\{N = k\}\right) \qquad &\text{indépendance de } N \text{ avec les } X_i
    \end{align*}
    Ensuite, écrivons~:
    \begin{align*}
        G_S(t) &= \sum_{n=0}^{+ \infty} \Bigg( \sum_{k=0}^{+ \infty}\proba\bigg(\bigg\{\sum_{i=1}^k X_i = n\bigg\}\bigg) \probaset{N = k} \Bigg)t^n\\
        &= \sum_{k=0}^{+\infty} \probaset{N=k} \sum_{n=0}^{+\infty} \proba \bigg( \bigg\{ \sum_{i=1}^k X_i = n \bigg\}\bigg) t^n \qquad &\text{théorème de Fubini positif}\\
        &= \sum_{k=0}^{+\infty} \probaset{N = k} G_{\sum_{i = 1}^k X_i} (t)\\
        &= \sum_{k=0}^{+\infty} \probaset{N = k} \prod_{i = 1}^k G_{X_i}(t) \qquad & \text{indépendance des } X_i\\
        &= \sum_{k=0}^{+\infty} \probaset{N=k} G_{X_1}(t)^k \qquad &\text{les } X_i \text{ sont identiquement distribuées}\\
        &= G_N\Big(G_{X_1}(t)\B ig)
    \end{align*}
    On retrouve bien le résultat, via un point de vue probabiliste.

\end{demonstration}


\section*{Questions de cours et remarques}

\subsection*{Questions de cours}
\begin{enumeratebf}
    \item lois discrètes usuelles (de Bernoulli, binomiale, géométrique, de Poisson)
    \item lois de densité usuelles
    \item inégalité de Markov
    \item inégalité de Bienaymé-Tchebychev
    \item covariance de la somme de variables aléatoires (et sa démonstration)
\end{enumeratebf}

\subsection*{Remarques}

\begin{itemize}
    \item Les énoncés sont à énoncer rigoureusement, en toute connaissance des hypothèses associées.
\end{itemize}

\end{document}